\clearpage
\section*{Supplementary material}

\setcounter{figure}{0}
\renewcommand{\thefigure}{S\arabic{figure}}
\renewcommand{\theHfigure}{S\arabic{figure}}

\begin{figure}[H]
    \centering
    \includegraphics[width=\textwidth]{plasticity_comparison.png}
    \caption{\textbf{Seed-paired comparison of instructive-driven and
    error-driven plasticity.} Each grey line joins instructive-driven and
    error-driven values obtained from the
    same root seed; colored points are individual seed-level values and black
    points with error bars show mean $\pm$ SEM across 20 seeds.
    \textbf{(Upper left)} Readout-compatibility effect, defined as aligned minus
    fixed-permutation output--target cosine in Fig.~\ref{fig:compatibility}A.
    \textbf{(Upper right)} Cue-exchange effect, defined as tuning similarity at
    exchange boundaries minus similarity at the same boundaries in the
    no-swap control in Fig.~\ref{fig:compatibility}B.
    \textbf{(Lower left)} Cue identity accuracy after 90\% LEC-like degradation
    in Fig.~\ref{fig:degradation}A.
    \textbf{(Lower right)} Nearest-position accuracy after 90\% MEC-like
    degradation in Fig.~\ref{fig:degradation}C. Both rules shared inputs,
    pretrained autoencoders, CA3 wiring seeds, storage order, corruption masks,
    and hyperparameters; only
    the CA3--CA1 update differed. Parameters were not optimized separately for
    each rule.}
    \label{fig:rule-comparison}
\end{figure}

\begin{figure}[p]
    \centering
    \begin{minipage}[t]{0.48\textwidth}
    \textbf{A}\par
    \includegraphics[width=\linewidth]{remapping_old_new_base.png}
    \end{minipage}
    \hfill
    \begin{minipage}[t]{0.48\textwidth}
    \textbf{B}\par
    \includegraphics[width=\linewidth]{remapping_old_new_err2.png}
    \end{minipage}

    \vspace{0.5em}

    \begin{minipage}[t]{0.48\textwidth}
    \textbf{C}\par
    \includegraphics[width=\linewidth]{remapping_endpoint_base.png}
    \end{minipage}
    \hfill
    \begin{minipage}[t]{0.48\textwidth}
    \textbf{D}\par
    \includegraphics[width=\linewidth]{remapping_endpoint_err2.png}
    \end{minipage}
    \caption{\textbf{Reference-map dynamics and heterogeneity after cue
    exchange.} \textbf{(A,B)} Similarity of the current CA1 tuning map to the
    final pre-exchange map (purple) and to the final map of the ensuing
    ten-lap block (orange), for instructive-driven (A) and error-driven (B)
    plasticity. Relative lap zero is the first lap learned after cue exchange.
    The adapted-map reference is retrospective and therefore equals one at the
    block endpoint by construction; the informative result is the abrupt
    crossover at exchange followed by smaller within-block changes.
    \textbf{(C,D)} Seed-averaged empirical cumulative distributions of
    unit-level tuning similarity between the pre-exchange and block-end maps
    for instructive-driven (C) and error-driven (D) plasticity. Purple denotes
    cue exchange and grey the matched no-swap control. For every curve, events
    and units were aggregated within each root seed before the mean and SEM
    were calculated across 20 seeds. These distributions describe population
    heterogeneity and do not treat units as independent replicates. All panels
    use the exact same EC--CA3 projection in the exchange and no-swap
    schedules.}
    \label{fig:remapping-dynamics}
\end{figure}
